% 列表间距与对齐
\usepackage{enumitem}
\let\tightlist\relax
\setlist{itemsep=0.5em, parsep=0.3em}

% 表格列宽
\setlength{\tabcolsep}{10pt}

% 图像绘制
\usepackage{pgfplots}
\pgfplotsset{compat=1.18}
\usetikzlibrary{positioning, trees, arrows.meta}

%%% 字体说明 %%%
% 主字体 → Noto Serif CJK SC（由 export.py 的 -V 参数传给 pandoc）
% 数学字体 → STIX Math（覆盖绝大多数数学符号）
% 不常见 Unicode 符号（如 ⟷⇏⇎ 等）→ 由 tools/unicode2latex.lua 提前替换为 LaTeX 命令
% XeLaTeX 不支持 RawFeature=fallback（那是 LuaLaTeX 的 luaotfload 特性）
% 故不在此处做字体级回退，而是靠 Lua filter 在标记→LaTeX 阶段完成转换

%%% 表格交替行底色 %%%
\usepackage[table]{xcolor}
\rowcolors{2}{gray!12}{white}

% 目录独立页码（罗马数字），目录结束后推到奇数页开始正文
\usepackage{etoolbox}
\AtBeginDocument{%
  \pretocmd{\tableofcontents}{\pagenumbering{Roman}}{}{}%
  \apptocmd{\tableofcontents}{%
    \newpage
    \ifodd\value{page}\relax\else\thispagestyle{empty}\null\newpage\fi
    \pagenumbering{arabic}%
  }{}{}%
}

% 页眉页脚：边缘标识（双面打印，主体2cm，标识突出到1cm）
\usepackage{fancyhdr}
\pagestyle{fancy}
\fancyhf{}

% 页眉向两侧延伸1cm（主体2cm → 标识位于1cm处）
\fancyhfoffset[RO]{1cm}
\fancyhfoffset[LE]{1cm}

% 页眉：外侧显示当前章节（灰底）
% \rightmark 取本页第一个标题，即"翻到这一页时所在的章节"；
% 三级标题都写 \markboth（left 与 right 同步），否则页面从 \section 起头时
% right 为空，页眉会变成一个空灰块。
\fancyhead[RO]{\colorbox{gray!15}{\rightmark}}
\fancyhead[LE]{\colorbox{gray!15}{\rightmark}}

% 页脚：外侧显示页码 + 内侧导出时间
\fancyfoot[RO]{\thepage}
\fancyfoot[LE]{\thepage}
\fancyfoot[LO]{EXPORTTIME}
\fancyfoot[RE]{EXPORTTIME}

%  plain 风格继承 fancy（解决 \section 首页丢失页眉）
\fancypagestyle{plain}{}
\AtBeginDocument{\pagestyle{fancy}}

\renewcommand{\headrulewidth}{0pt}
\renewcommand{\footrulewidth}{0pt}
\makeatletter
\renewcommand{\sectionmark}[1]{\markboth{#1}{#1}}
\renewcommand{\subsectionmark}[1]{\markboth{#1}{#1}}
\renewcommand{\subsubsectionmark}[1]{\markboth{#1}{#1}}
\makeatother
